\begin{textsample}{POS Dim 1 – 2015\_c5 – Score 10.00 – 2015\_c5/enade\_2015\_tec\_gest\_publ\_0d4a20dfe1.txt}  \label{ex:f1_pos_095}
CAPÍTULO 5 DISTRIBUIÇÃO DOS CONCEITOS 5.1 \textbf{PANORAMA} NACIONAL DA DISTRIBUIÇÃO DOS CONCEITOS O Conceito Enade é um indicador de qualidade que avalia o desempenho dos estudantes a partir dos resultados do Enade.
Os valores possíveis vão de 1 ( pior situação ) a 5 ( melhor situação ).
As Unidades de Observação com apenas um ou sem nenhum concluinte participante não obtêm o Conceito Enade, ficando Sem Conceito ( SC ).
A Tabela 5.1 apresenta a quantidade e distribuição de cursos de Tecnologia em Gestão Pública participantes do Enade/2015, por faixa de conceito e Grande Região.
Mais uma vez, ressaltando que a diferença entre os cursos \textbf{tabulados} neste capítulo e no capítulo 2 corresponde aos cursos sem conceito, em princípio, aqueles sem alunos concluintes que participassem da prova ou que tivessem tão somente um aluno concluinte18.
Observando -se os dados da Tabela 5.1, nota -se que, dos 65 cursos participantes, 24 ( 36,9 \% ) classificaram -se com conceito 3, o valor modal.
Este foi também o conceito modal na região Sul ( 62,5 \% ).
Na região Centro-Oeste, a distribuição dos cursos entre os conceitos exibiu um comportamento \textbf{bimodal} : cinco cursos foram alocados no conceito 2 e outros cinco obtiveram conceito 4 ( 33,3 \% cada ).
O conceito 4 foi o segundo mais frequente em nível nacional ( 29,2 \%, correspondendo a 19 cursos ) e conceito 2, o terceiro ( 20,0 \%, correspondendo a 13 cursos ).
Houve, ainda, dois cursos ( 3,1 \% ) que receberam conceito 5 e três ( 4,6 \% ) cursos que receberam conceito 1.
Dos 65 cursos de Tecnologia em Gestão Pública, quatro cursos ( 6,2 \% ) ficaram sem conceito ( SC ). 18 Estes cursos com somente um concluinte não foram considerados no capítulo 2.
A região Norte participou com quatro cursos ou 6,2 \% do total nacional.
Nesta região, o conceito 4 foi atribuído a dois cursos dos quatro participantes, o que \textbf{equivale} a 50,0 \% do total regional, além de ser o conceito modal na região.
Nenhum curso ficou sem conceito ( SC ) ou com conceitos 2 ou 5.
Dois cursos foram avaliados com conceitos 1 e 3 ( um curso, cada ), correspondendo a 25,0 \%, cada.
A região Nordeste participou com 11 cursos ou 16,9 \% do total nacional.
Nesta região, quatro cursos, 36,4 \% em termos regionais, obtiveram conceito 4, o conceito modal para a região.
Os conceitos 2 e 3 foram atribuídos a, respectivamente, dois e três cursos ( 18,2 \% e 27,3 \% ).
Um curso obteve o conceito 1, e um outro ficou sem conceito ( SC ), correspondendo a 9,1 \% cada.
Nessa região, nenhum curso recebeu conceito 5.
Dos 19 cursos participantes da região Sudeste ( 29,2 \% do total nacional ), o conceito 4 também foi o conceito modal, tendo sido obtido por sete cursos ( 36,8 \% ).
O conceito 1 não foi atribuído a qualquer curso, e os conceitos 2 e 3, a dois e seis cursos, correspondendo a 10,5 \% e 31,6 \%, respectivamente.
O conceito 5 foi atribuído a dois cursos ( 10,5 \% ).
E outros dois cursos ( 10,5 \% ) ficaram sem conceito ( SC ).
Os 16 cursos da região Sul ficaram distribuídos em quase todas as faixas de conceitos.
Há predominância do conceito 3 ( conceito modal ), atribuído a dez cursos ( 62,5 \% ), ou seja, dez dos 16 cursos participantes na região Sul.
O conceito 2 foi atribuído a quatro cursos ( 25,0 \% ) e o conceito 4, a somente um curso ( 6,3 \% ).
Apenas um curso na região Sul ficou sem conceito ( SC ), enquanto nenhum recebeu os conceitos 1 ou 5.
Na região Centro-Oeste foi 15 o número de cursos participantes.
Destes, cinco receberam conceito 2 e outros cinco foram alocados ao conceito 4 ( 33,3 \% cada ).
Nesta região, nenhum curso ficou sem conceito ( SC ) ou obteve conceito 5.
Já quatro cursos receberam conceito 3, correspondendo a 26,7 \% do total regional, e um curso foi avaliado com conceito 1 ( 6,7 \% ).
O Gráfico 5.1 apresenta a distribuição do Conceito Enade segundo Grande Região : \textbf{poligonais} mais à esquerda representam Grandes Regiões cuja distribuição apresenta os menores conceitos, e \textbf{poligonais} mais à direita, Grandes Regiões com melhores conceitos.
A região Sudeste ( linha vermelha ) apresenta as distribuições com valores maiores.
Os cursos em IES das demais regiões, Norte ( linha azul ), Nordeste ( linha verde ), Sul ( linha \textbf{roxa} ) e Centro-Oeste ( linha preta ) apresentam uma situação pior. 5.2 CONCEITOS POR CATEGORIA ADMINISTRATIVA E POR GRANDE REGIÃO A Tabela 5.2 apresenta a distribuição dos cursos participantes do Enade/2015 de Tecnologia em Gestão Pública, por Categoria Administrativa, de acordo com os conceitos por eles alcançados, segundo as Grandes Regiões brasileiras.
Dos 65 cursos participantes, 17 ( 26,2 \% ) eram \textbf{ministrados} em Instituições \textbf{Públicas} e 48 ( 73,8 \% ), em Instituições Privadas.
De acordo com as informações da Tabela 5.2, em termos nacionais, entre Instituições \textbf{Públicas} e Privadas observa -se que todas as duas IES que receberam o conceito 5 eram Privadas.
Dos 17 cursos participantes de IES Públicas, o conceito 3 foi o valor modal, atribuído a dez cursos ( 58,8 \% ).
Entre os demais cursos de Instituições \textbf{Públicas} participantes, sete cursos ( 41,2 \% ) foram avaliados com conceito 4.
Nesta categoria, nenhum curso ficou sem conceito ( SC ) ou recebeu os demais conceitos.
Na rede privada, também, o conceito modal foi 3, atribuído a 14 cursos dos 48 da categoria.
Entre os demais cursos participantes, três receberam conceito 1, e 13, conceito 2.
O conceito 4 foi atribuído a 12 cursos e o conceito 5, a dois cursos.
Nesta Categoria Administrativa, outros quatro cursos ficaram sem conceito ( SC ).
Na análise por região, observa -se que, na região Norte, as Instituições Privadas participaram com somente um curso ( 25,0 \% do total regional ), a menor das proporções desta categoria comparada as demais regiões brasileiras, e este foi alocado no conceito 1.
Nenhum curso ficou sem conceito ( SC ) ou recebeu os conceitos \textbf{restantes} nessa categoria e região. categoria e região.
As Instituições Públicas participaram com três cursos na região Norte ( 75,0 \% do total regional ), a maior das proporções desta categoria em relação as demais regiões brasileiras, dos quais dois obtiveram conceitos 4, o valor modal, seguido por um curso ( 33,3 \% ) que recebeu conceito 3.
Nenhum curso ficou sem conceito ( SC ) ou recebeu os demais conceitos nesta combinação de categoria e região.
Na região Nordeste, a rede privada concentrou seis dos 11 cursos participantes, o equivalente a 54,5 \% do total da região.
Dentre os cursos oferecidos por IES Privadas no Nordeste, um curso ficou alocado no conceito 1, outro curso obteve conceito 3, mais um curso recebeu conceito 4 e um último curso ficou sem conceito ( SC ).
Além destes, dois cursos obtiveram conceito 2, conceito modal, nessa categoria, e nenhum obteve conceito 5.
As Instituições \textbf{Públicas} dessa região participaram com cinco cursos ( 45,5 \% do total regional ), dos quais três ( 60,0 \% ) obtiveram conceito 4, o valor modal, seguido por dois cursos que receberam conceito 3 ( 40,0 \% ).
Nenhum curso ficou sem conceito ( SC ) ou recebeu os demais conceitos nesta combinação de categoria e região.
Na região Sudeste, a proporção de cursos da rede privada foi de 89,5 \%, correspondendo a 17 dos 19 cursos participantes.
Nas IES Privadas da região Sudeste, o conceito modal foi 4 ( seis cursos ).
Nesta combinação de Categoria Administrativa e Grande Região, dois cursos de IES Privada não receberam conceito ( SC ), enquanto nenhum curso recebeu conceito 1.
Os demais foram avaliados com conceito 2 ( dois cursos ), conceito 3 ( cinco cursos ) e conceito 5 ( dois cursos ).
Entre os dois cursos oferecidos por Instituições \textbf{Públicas} na região Sudeste, um curso recebeu o conceito 3 e outro foi alocado no conceito 4.
Na região Sudeste, nesta categoria, nenhum curso ficou sem conceito ( SC ) ou recebeu os conceitos \textbf{restantes}.
As Instituições Privadas concentraram dez dos 16 cursos participantes da região Sul, 62,5 \% do total regional.
Desses, oito ficaram com conceitos 2 e 3 ( quatro cursos cada ).
Um único curso foi associado ao conceito 4, enquanto um curso ficou sem conceito ( SC ) e nenhum recebeu os conceitos 1 ou 5.
As Instituições \textbf{Públicas} na região Sul participaram com seis cursos ( 37,5 \% ), e todos eles foram avaliados com conceito 3.
Nenhum curso ficou sem conceito ( SC ) ou recebeu os demais conceitos nessa categoria e região.
Na região Centro-Oeste, 14 dos 15 cursos participantes eram de Instituições Privadas ( 93,3 \% em termos regionais ), a proporção mais elevada desta categoria das regiões brasileiras.
Destes, cinco concentraram -se no conceito 2, conceito modal.
Dentre os demais, apenas um curso obteve conceito 1, quatro receberam conceito 3, e outros quatro foram alocados no conceito 4.
Nenhum curso ficou sem conceito ( SC ) ou recebeu conceito 5 nesta categoria.
Somente um curso foi oferecido por Instituições \textbf{Públicas}, e este recebeu conceito 4.
Nesta região, nenhum curso de IES \textbf{pública} ficou sem conceito ( SC ) ou recebeu os conceitos \textbf{restantes} dessa categoria e região.
O Gráfico 5.2 apresenta a distribuição do Conceito Enade segundo a Categoria Administrativa da IES.
Os cursos mantidos por IES Públicas ( linha azul ) apresentam uma \textbf{poligonal} mais à direita do que os oferecidos por IES Privadas ( linha verde ) e, consequentemente, uma distribuição de Conceitos Enade com valores maiores. 5.3 CONCEITOS POR ORGANIZAÇÃO ACADÊMICA E POR GRANDE REGIÃO Na Tabela 5.3, encontra -se a distribuição dos conceitos atribuídos aos cursos participantes do Enade/2015 na Área de Tecnologia em Gestão Pública, por Organização Acadêmica, segundo as Grandes Regiões brasileiras.
Dos 65 cursos de Tecnologia em Gestão Pública participantes, 36 eram oferecidos em Universidades, 16 em Centros Universitários e 13 em Faculdades.
Esta distribuição corresponde a, respectivamente, 55,4 \%, 24,6 \% e 20,0 \% dos cursos.
De acordo com os dados apresentados, os dois cursos avaliados com conceito 5 não eram vinculados a Universidades.
Dois cursos oferecidos em Universidades ficaram sem conceito ( SC ).
Esse tipo de Organização Acadêmica teve o conceito 3 como modal, atribuído a 16 cursos ( 44,4 \% ).
Os demais cursos de Universidades avaliados receberam os conceitos 2 ( sete cursos ) e 4 ( 11 cursos ).
Nenhum curso recebeu conceito 1 ou, como já mencionado, conceito 5.
Entre os cursos em Centros Universitários, o conceito modal foi 4, atribuído a seis cursos ( 37,5 \% ).
Neste tipo de Organização Acadêmica, somente um curso ficou sem conceito ( SC ), enquanto um outro obteve conceito 1.
Os outros cursos receberam os conceitos 2 ( três cursos ) e 3 ( cinco cursos ).
Nenhum curso recebeu conceito 5.
Nas Faculdades, um dos 13 cursos ficou sem conceito ( SC ), e seis cursos ficaram alocados nos conceitos 2 e 3 ( três cursos em cada ), seguidos de seis cursos que obtiveram conceitos 1, 4 e 5 ( um desses conceitos para cada dois cursos ).
Considerando -se separadamente as regiões brasileiras, verifica -se que, na região Norte, as Universidades concentraram três dos quatro cursos participantes.
Dois cursos ficaram alocados ao conceito 4, conceito modal, seguidos de um curso que obteve conceito 3.
Nenhum curso ficou sem conceito ( SC ) ou recebeu os demais conceitos nesta categoria e região.
Os Centros Universitários da região Norte foram representados por somente um curso, e este recebeu o conceito 1.
Não houve presença de cursos entre as Faculdades.
Na região Nordeste, as Universidades participaram com oito dos 11 cursos na Área de Tecnologia em Gestão Pública.
Seis cursos obtiveram os conceitos 3 e 4 ( três cursos com cada um dos conceitos ) e um curso foi alocado no conceito 2, enquanto um outro ficou sem conceito ( SC ).
Nenhum dos cursos oferecidos por Universidades no Nordeste recebeu os conceitos 1 ou 5.
Os Centros Universitários \textbf{contaram} com somente um curso participante na região Nordeste, e este obteve conceito 4.
As Faculdades foram representadas por dois cursos na região Nordeste, dos quais um recebeu conceito 1 e o outro ficou com conceito 2.
Na região Sudeste, as Universidades concentraram dez dos 19 cursos da região.
Entre os cursos em Universidades, nessa região, oito cursos obtiveram os conceitos 3 e 4 ( quatro cursos com cada um dos conceitos ), enquanto dois cursos ficaram com conceito 2.
Nenhum curso ficou sem conceito ( SC ) ou recebeu os conceitos 1 ou 5 nesta combinação de categoria e região.
Os Centros Universitários participaram com quatro cursos na região Sudeste, dos quais dois obtiveram o conceito modal, 4.
Um curso ficou sem conceito ( SC ) e um outro curso recebeu conceito 3.
Nenhum curso recebeu os conceitos \textbf{restantes} nessa categoria e região.
As Faculdades foram representadas por cinco cursos na região Sudeste, dos quais dois receberam conceito 5 ( conceito modal ).
E dois cursos obtiveram os conceitos 3 e 4 ( um curso cada ).
Um curso ficou sem conceito ( SC ) enquanto nenhum recebeu os conceitos 1 ou 2 nessa categoria e região.
Dos 16 cursos da região Sul, 12 eram de Universidades, para os quais o conceito modal foi 3, atribuídos a oito cursos, seguidos de três cursos que obtiveram conceito 2.
Nesse tipo de organização, um dos cursos ficou sem conceito ( SC ) enquanto nenhum recebeu os demais conceitos nesta região.
Dos três cursos participantes de Centros Universitários da região Sul, dois obtiveram conceito 3, o conceito modal, seguidos de um curso, que foi alocado no conceito 2.
Nenhum curso ficou sem conceito ( SC ) ou recebeu os conceitos \textbf{restantes} nessa categoria e região.
Somente um curso foi vinculado a Faculdades na região Sul, e este recebeu conceito 4, enquanto nenhum ficou sem conceito ( SC ) ou recebeu os demais conceitos nessa categoria e região.
Na região Centro-Oeste, três dos 15 cursos eram de Universidades.
Nesse tipo de organização, nenhum curso ficou sem conceito ( SC ) ou com os conceitos 1, 3 ou 5.
O conceito modal foi 4, atribuído a dois cursos, seguidos de um curso que obteve conceito 2.
Os Centros Universitários da região Centro-Oeste \textbf{contaram} com sete cursos, que obtiveram conceitos 2 ( dois cursos ), 3 ( dois cursos ) e 4 ( três cursos, o conceito modal ).
Nenhum curso ficou sem conceito ( SC ) ou recebeu os conceitos 1 ou 5 na categoria.
Dos cinco cursos oferecidos em Faculdades na região Centro-Oeste, quatro ficaram alocados nos conceitos 2 e 3 ( dois cursos em cada conceito ), e um curso obteve conceito 1.
Nenhum curso ficou sem conceito ( SC ) ou recebeu os conceitos 4 ou 5 nesta categoria e região.
O Gráfico 5.3 apresenta a distribuição do Conceito Enade segundo a Organização Acadêmica da IES.
Os cursos em Universidades ( linha azul ) e Centros Universitários ( linha verde ) apresentam \textbf{poligonais} bem próximas, indicando semelhança entre as distribuições dos Conceitos Enade.
Os cursos oferecidos em Faculdades ( linha vermelha ) apresentam uma \textbf{poligonal} \textbf{cruzando} as outras duas e, consequentemente, uma distribuição com maior variância de Conceitos Enade.

% matched lemmas: bimodal, contar, cruzar, equivaler, ministrar, panorama, poligonal, pública, restante, roxa, tabular
\end{textsample}
